\section[第10章\q 还想要更多吗？\q 这里有一些要做的事情]{第10章\\还想要更多吗？\\这里有一些要做的事情}

许多人可能并不满足于我说的这10个步骤，他们认为这些步骤更是一种思想而不是实际行动。而我认为，理解这一思想和行动同样重要。有许多人愿意去做而不愿意思考，也有许多人愿意思考而不愿意付诸行动。而我既愿意思考也愿意行动。我喜欢新思想，也乐于付诸行动。

那么，对于那些想行动起来的人来说，应该如何开始呢？在此我想简要地介绍一下我是怎样做的。

停下你手头的活儿。换句话说，就是先停下来，评估一下你的做法中哪些有效，哪些无效。做同一件事情却希望有不同的结果是神志不清的表现。不要做那些无效的事情，找一些有效的事情去做。

寻找新的思想。为了寻找投资的灵感，我经常到书店搜寻提供独特的、与众不同的主张的书，我把它们称为模式。我买那些介绍各种“模式”的书，这些模式是我以前不知道的。例如，在书店，我找到了乔尔·莫斯科维茨的《收益率达到16％的方法》，我买下了这本书并一口气读完。

开始行动！

在第二周的星期四，我开始按照书上说的话一步步行动。我去律师事务所和银行寻找廉价房地产交易的机会。大部分人并不采取行动，或者被别人说服不去应用所学到的任何新模式。我的邻居就曾经对我说，16％的收益率是不可能实现的，但我没有听他的，因为他从来没有试过。

找一些做过你想做的事情的人，请他们和你一起共进午餐，向他们请教一些诀窍和技巧。就拿16％的税收留置权来说吧，我到税务办公室同一位工作人员见面，我发现她也在做税收留置权投资，于是立即邀请她一起吃午饭。她很兴奋地告诉我她所知道的有关这种投资的所有方法。甚至在吃过饭后，她又用了整整一下午来向我说明投资的全过程。

到了第二天，我就在她的帮助下，找到了两笔大买卖，从此我每年就能获得16％的利息收益。我花了一天来读有关的书，用一天来采取行动，用一个小时吃午饭，又用一天找到了两笔大买卖。

参加辅导班并购买相关磁带。我在报纸上寻找让我感兴趣的辅导班的广告，有许多是免费的，也有一些只收取一小笔费用。我也参加一些很贵的研讨班，因为这些地方所讨论的内容正是我想学的。就是因为学习了这些课程，我才会变得富有，不用出去辛苦地工作。我有许多朋友从不参加这种辅导班，并说我是在浪费钱，然而他们至今仍在干着和以前一样的工作。

提出多份报价。如果我想要买一处房产，我会看多处房产并给出一个一般的报价。如果你不知道什么是正确的报价，那我也不知道，这是房地产中介的工作，要由它们来提出报价。我尽可能少做这方面的工作。

一位朋友希望我告诉她如何购买公寓。所以在一个周六，她、房地产中介和我一起去看了6处公寓。其中4处不太好，另外两处不错。我提议对所有6处都报价，价格为卖主要价的一半。她和中介非常吃惊，认为这样做未免太无礼了，恐怕会冒犯卖主们。但是我觉得，房地产中介并不那么尽力，所以他们常常什么也不做而是去寻找一处更好的房产。

后来他们一个报价也没报，而那个中介仍然在寻找一笔价格“合适”的交易。其实，你根本就不知道什么价格才是“合适”的，除非有另一处同样的交易作为参照。大部分卖主的要价过高，很少有要价低于标的物的实际价值。

这个故事的中心是：多发出几份报价。没卖过东西的人，对想卖出东西的迫切心情是不会理解的。我有一处房产，想在数月之内卖掉，当时我愿意接受任何报价，不会在意价格有多低，即使他们只给我10头猪我也会非常高兴。报价本身并不重要，关键是有人感兴趣。也许我会反过来建议对方以一个猪圈交换房产也不一定呢，游戏就是这样的。记住，做买卖就是一场有趣的游戏。你提出报价之后，可能对方就会说：“成交。”

我还经常使用“免责条款”来做报价。例如在房地产交易协议上，我会加上一条“须得到我的商业伙伴的同意”。我从不指明我的商业伙伴到底是谁，大部分人都不知道我的商业伙伴其实就是我的小猫。如果他们接受我的报价，而我又想反悔的话，我就给家里的小猫打电话。我讲这个荒唐故事的目的就是为了说明，买卖游戏简单得简直让人难以置信。

所以，我觉得许多人太过严肃了，反而把事情弄得太复杂。

寻找一桩好生意、一家好企业、一位合适的人、一位合适的投资者，或任何类似的东西，就如同约会一样。你必须到市场上去和许多人谈，做许多报价、还价、谈判、拒绝或者接受。我知道有些单身的人宁可在家里坐等电话铃响，但是，除非你是辛迪·克劳馥或者汤姆·克鲁斯，否则你最好还是到市场上去，即使只是一家超市也好。从寻找、报价、拒绝、谈判到成交，几乎是人的一生中要经历的全部过程。

每月在某一地区慢跑、散步或驾车10分钟左右。在慢跑的过程中，我曾发现我的投资中最好的房地产项目。有一年，我常常固定在某一邻近的地区慢跑，为的是发现某些变化。一桩交易要赢利，必须具备两个条件：一是廉价，二是有变化。市场上有许多廉价交易，但只有存在变化，才能使廉价交易变得有利可图。因此，当我慢跑的时候，经过的地方都是我想要投资的区域。通过反复观察，我注意到一些细微的变化。

我会注意到有的出售招牌挂了很长时间，那意味着卖主很可能急于成交。看到卡车进进出出，我会停下来和司机交谈。我还与邮政货车的司机聊天，从这些人口中可以得到有关某一地区详细得惊人的信息。

我找到一个市场行情很差的地区，是那种人人唯恐避之不及的地区。在一年里我不时开车来到这里，以观察有没有好转的迹象。我和零售店主，特别是那些新搬来的交谈，以弄清楚他们搬来的原因。这样做每月只需要花很少时间，同时我还能做其他的事情，比如锻炼锻炼身体，或去商店走走看看等。

至于股票，我喜欢彼得·林奇写的《称雄华尔街》一书中介绍的选择有升值潜力股的方法。我发现寻找有升值潜力的事物的方法都是相同的，不管他们是房地产、股票、共同基金、新公司、新宠物、新房子、新配偶还是一包洗衣粉。

过程往往是一致的。你要知道你在寻找什么，然后开始行动！

为什么消费者总是穷人？每当超市大减价，如卫生纸打折时，消费者就会涌入超市抢购。而当股票市场上出现股价下跌时，也就是大多数人所说的股市下挫或回调时，购买者却急于从中逃出。当超市商品涨价时，人们转到其他商店购物。而当股市上升时，购买者却大举买入股票。

关注适当的地方。一位邻居以10万美元购买了一套公寓，我则以5万美元购买了与之相邻的相同的一套公寓。他告诉我他在等着价格上涨。我对他说他能获得的利润在他购买这套公寓时就已经确定了，而并不是在他出售时确定的。他是通过一位房地产经纪人来进行交易的，而这位经纪人却并没有属于自己的房产。我是在一家银行的破产清偿部购买的。我花了500美元参加了一个学习班，学习如何在破产清偿部购买这类房产。我的邻居认为花500美元上一个学习班未免太贵了，他说他没有钱，也没有时间，所以他只能寄希望于价格上升。

我首先寻找想买进的人，然后才去找想卖出的人。一位朋友想买一片地，他有钱，但没时间去找。我发现了一处地产，比我朋友想要的面积大一些。我打电话给朋友，他说想要其中一片，于是我买下了土地，然后把那一片卖给了他。我一直没有卖掉剩余的土地，它一直为我带来收益。这个故事的中心是：买下馅饼并把它切成小块。大部分人寻找的只是自己能付得起的东西，所以他们只看到较小的东西。他们只买一小块馅饼，却总是付出更高的价格。只盯着小生意的人是不会有大突破的。你如果想致富，就要首先考虑较大的生意。

零售商喜欢提供数量折扣，就是因为大部分商人都喜欢大额购买的人。所以即使你的投资规模很小，你也可以多考虑考虑大生意。当我的公司想购买电脑时，我就打电话给几位朋友，问他们是否也要买电脑。接着我们到不同的零售商那里进行谈判，尽量压低价格，因为我们购买电脑的数量很大。我也以同样的方式买卖股票。小规模投资人善于小规模的动作，因为他们思考的范围太狭窄，他们总是单干，从不协同作战。

学习前人经验。所有上市的大公司都是从小公司起家的。桑德斯上校在60多岁失去了所有财产之后才致富。比尔·盖茨在30岁以前就成了世界上最富有的人之一。

行动的人总会击败不行动的人。

以上这些只是我过去曾经做过的事情中的一小部分，我将继续发现投资机会。最重要的是“做过”和“去做”。在书中我曾多次提到，在你得到财富的回报以前必须先行动。那么现在就行动吧！
